\documentclass[11pt]{article}

\usepackage[a4paper,margin=1in]{geometry}
\usepackage{amsmath,amssymb,mathtools}
\usepackage{hyperref}

\newcommand{\ones}{\mathbf{1}}
\newcommand{\diag}{\operatorname{diag}}
\newcommand{\KL}{\operatorname{KL}}
\newcommand{\Tset}{\mathcal{T}}

\title{Perspectives on Learning Reverse Dynamics in Discrete Diffusion}
\author{}
\date{}

\begin{document}
\maketitle

\section{Setup and the structural gap with Euclidean diffusion}

Let $\mathcal X=\{1,\dots,n\}$ and $(h^t)_{t=0}^T$ be histograms on $\mathcal X$:
\[
h^{t+1}=P^\top h^t,\qquad t=0,\dots,T-1.
\]
Reverse kernels are matrices $Q^t\in\mathbb R_+^{n\times n}$ acting by
\[
h^t=(Q^t)^\top h^{t+1}.
\]
Admissible kernels are
\[
\Tset(h^{t+1},h^t)
=\{K\ge 0:\ K\ones=\ones,\ K^\top h^{t+1}=h^t\}.
\]

The key distinction with Euclidean diffusion is:
\[
\text{discrete reverse model: }(i,t)\mapsto Q^t(i,\cdot)\in\Delta_n,
\qquad
\text{continuous reverse model: }(x,t)\mapsto v_t(x)\in\mathbb R^d.
\]
So discrete learning is distribution-valued, while score/velocity learning is vector-valued. This structural mismatch is the origin of the sampling-cost gap.

\section{Direct conditional regression: precise objective and limitations}

\subsection*{Population objective decomposition}
Let $Q_\star^t(i,\cdot)$ be the true reverse conditional and $\mu_t(i)=h_i^{t+1}$ the context law. Define
\[
\mathcal R(\theta)
=\frac1T\sum_{t=0}^{T-1}\sum_{i=1}^n \mu_t(i)\,
\mathbb E_{j\sim Q_\star^t(i,\cdot)}[-\log Q_\theta^t(i,j)].
\]
Then
\[
\mathcal R(\theta)
=C_\star+\frac1T\sum_{t=0}^{T-1}\sum_{i=1}^n \mu_t(i)\,
\KL\!\big(Q_\star^t(i,\cdot)\,\|\,Q_\theta^t(i,\cdot)\big),
\]
where $C_\star$ does not depend on $\theta$. Hence NLL is statistically aligned with conditional recovery.

\subsection*{Admissibility is not automatic}
Even if each row of $Q_\theta^t$ is normalized (softmax output),
\[
Q_\theta^t\ones=\ones
\]
does \emph{not} imply
\[
(Q_\theta^t)^\top h^{t+1}=h^t.
\]
A precise discrepancy metric is
\[
\Delta(\theta)
\coloneqq
\frac1T\sum_{t=0}^{T-1}
\left\|(Q_\theta^t)^\top h^{t+1}-h^t\right\|_1.
\]
One can train with
\[
\min_\theta\ \mathcal R(\theta)+\lambda\,\Delta(\theta),
\]
or use explicit projection/parameterization so $Q_\theta^t\in\Tset(h^{t+1},h^t)$ by construction.

\section{Sampling complexity: where brute-force fails}

If $\mathcal X$ has size $n$, sampling one backward trajectory from dense kernels costs roughly
\[
\mathcal O(T\,n)
\]
without preprocessing (or $\mathcal O(T)$ draws after $\mathcal O(Tn^2)$ alias preprocessing).
For very large $n$ (or infinite discrete spaces), dense categorical outputs become the bottleneck.

Structured factorizations reduce this:
\begin{itemize}
\item autoregressive tokenization: replace one $n$-way draw by a product of low-arity draws;
\item local-move kernels with degree $d\ll n$: per-step cost $\mathcal O(d)$.
\end{itemize}
This is a primary practical reason to prefer structured reverse parameterizations over unconstrained full-conditionals.

\section{Binary Langevin viewpoint and diffusion-compatible extension}

For $\mathcal X=\{0,1\}^d$, denote by $x^{(k)}$ the bit-flip of coordinate $k$.
Langevin-type binary samplers (as in arXiv:2502.00557) use gradient-informed local moves instead of full-state categoricals.
This suggests the following diffusion-compatible family.

\subsection*{Sparse reverse kernel parameterization}
Choose neighbor set
\[
\mathcal N(x)=\{x^{(1)},\dots,x^{(d)}\}.
\]
Define logits $s_\theta^t(x,k)$ and local probabilities
\[
\pi_\theta^t(k\mid x)
=\frac{\exp(s_\theta^t(x,k))}{\sum_{\ell=1}^d \exp(s_\theta^t(x,\ell))}.
\]
Set a sparse kernel
\[
Q_\theta^t(x,y)=
\begin{cases}
\alpha_t\,\pi_\theta^t(k\mid x), & y=x^{(k)},\\
1-\alpha_t, & y=x,\\
0, & \text{otherwise},
\end{cases}
\]
with $\alpha_t\in(0,1]$.
This guarantees row-stochasticity and yields $\mathcal O(d)$ sampling.

\subsection*{Connection with score-like quantities}
If one has a score surrogate $g_t(x)\approx \nabla_x \log h^t(x)$ (or discrete differences), one can set
\[
s_\theta^t(x,k)\approx \beta_t\,\delta_k \log h^t(x),
\]
where $\delta_k$ is a discrete directional difference. This mimics Langevin-style local preference while remaining in a Markov-kernel formulation.

\subsection*{Diffusion constraints}
The target constraints remain exactly
\[
h^{t+1}=P^\top h^t,\qquad
h^t=(Q^t)^\top h^{t+1},\qquad
Q^t\in\Tset(h^{t+1},h^t).
\]
Hence the concrete training problem is:
\[
\min_\theta\ \mathcal R(\theta)+\lambda\,\Delta(\theta)
\quad\text{with sparse }Q_\theta^t.
\]
This is the precise bridge: keep diffusion equations unchanged, replace dense conditionals by structured local kernels inspired by binary Langevin updates.

\section{Discrete optimal transport interpretation on the binary cube}

For each time step, set
\[
a=h^{t+1},\qquad b=h^t,\qquad \pi^t=\diag(a)\,Q^t.
\]
Then
\[
Q^t\in\Tset(a,b)
\quad\Longleftrightarrow\quad
\pi^t\in\Pi(a,b),
\]
where
\[
\Pi(a,b)=\{\pi\ge 0:\ \pi\ones=a,\ \pi^\top\ones=b\}.
\]
So reverse-kernel learning is equivalent to learning couplings with prescribed marginals.

On $\{0,1\}^d$, define a geometry-aware cost $c(x,y)$, for instance
\[
c(x,y)=\|x-y\|_0,
\]
or a strict local-move cost with $c(x,y)=+\infty$ when $y\notin \mathcal N(x)\cup\{x\}$.
Each reverse step can then be posed as
\[
\min_{\pi^t\in\Pi(h^{t+1},h^t)} \langle C,\pi^t\rangle,
\qquad C_{xy}=c(x,y),
\]
which is a discrete OT problem on the cube.

To integrate learning, use a reference coupling
\[
\bar\pi_\theta^t=\diag(h^{t+1})\,\bar Q_\theta^t
\]
from a score/Langevin-inspired sparse kernel, and solve
\[
\min_{\pi^t\in\Pi(h^{t+1},h^t)}
\ \langle C,\pi^t\rangle
+\varepsilon\,\KL(\pi^t\|\bar\pi_\theta^t).
\]
This yields a sharp decomposition:
\begin{itemize}
\item $C$ imposes transport geometry (Hamming/locality);
\item $\bar\pi_\theta^t$ injects learned directional prior;
\item $\Pi(h^{t+1},h^t)$ enforces diffusion-consistent marginals.
\end{itemize}
Hence the perspective is not only ``fit a conditional'' but ``fit a time-indexed OT coupling field with learned reference dynamics''.

\section{Sharper open problems}

\begin{itemize}
\item Characterize model classes where $Q_\theta^t\in\Tset(h^{t+1},h^t)$ holds by design with low sampling complexity.
\item Quantify the expressivity gap between dense kernels and sparse local kernels under fixed computational budget.
\item Establish when score-like local logits recover the Bayes reverse kernel (exactly or approximately) on $\{0,1\}^d$.
\item For sequence models, combine causality factorization and diffusion-time conditioning without serializing all dimensions.
\item Evaluate models on a triplet of criteria: NLL, admissibility gap $\Delta(\theta)$, and effective wall-clock sampling cost.
\end{itemize}

\end{document}
